\usepackage[noend,noline]{algorithm2e}
\usepackage{amsmath}
\usepackage[polish]{babel}
\usepackage{booktabs}
\usepackage{euler}
\usepackage{fontspec}
\usepackage{graphicx}
\usepackage[export]{adjustbox}
\usepackage{hyperref}
\usepackage{minted}
\usepackage{multicol}
\usepackage{pgfplots}
\pgfplotsset{compat=1.18}
\usepackage{tikz}
\usepackage{xcolor}

\def \LongInput#1{%
  \ifWydruk
    \begin{frame}{\FrameTitle}
      #1
    \end{frame}
  \else
    \input #1
  \fi
}

\def \IncludeDiagram#1#2{%
  \ifWydruk
    \includegraphics[#1]{./images/white/#2}
  \else
    \includegraphics[#1]{./images/black/#2}
  \fi
}

\usetheme{metropolis}

\ifWydruk
  \definecolor{myorange}{HTML}{005500}
\else
  \usecolortheme{owl}
  \definecolor{myorange}{HTML}{FFA500}
  \definecolor{jasnozolty}{RGB}{255, 255, 200}
  \setbeamercolor{normal text}{fg=jasnozolty}
\fi
\definecolor{myred}{HTML}{FF0000}
\definecolor{mygreen}{HTML}{33BB00}
\definecolor{myblue}{HTML}{7777FF}
\definecolor{mygray}{HTML}{BAB5C8}

\setbeamersize{text margin left=5mm}

% Jeśli to polecenie przestanie działać z~powodu zmiany Pythona 3 na
% nowszą wersję, należy skopiować pliki
% ~/.local/lib/python3.10/site-packages/pygments/styles/bw-with-orange-comments.py
% ~/.local/lib/python3.10/site-packages/pygments/styles/__init__.py
% na nowe miejsca.
\usemintedstyle[go]{bw-with-orange-comments}

\usetikzlibrary{automata,positioning,arrows,arrows.meta}

\tikzset{
  ->,
  >=stealth,
  initial text={\it start},
  every state/.style={inner sep=2pt,minimum size=3mm},
  double=black,
  double distance=2pt,
  node distance=15mm,
}

\frenchspacing
\def \g{\color{myorange}}
\let \ttt = \tt
\def \tt{\g\ttt}
\let \iit = \it
\def \it{\g\iit}

\author{dr~inż.~Marcin Ciura}
\institute{%
  Wydział Bezpieczeństwa i Informatyki\\
  Uniwersytet Komisji Edukacji Narodowej}
\date{}

\setsansfont{RecMonoCasual-Regular-1.085.ttf}[
  Path=/Users/marcin/Library/Fonts/,
  Contextuals=NoAlternate,
  BoldFont=RecMonoCasual-Bold-1.085.ttf,
  ItalicFont=RecMonoCasual-Italic-1.085.ttf,
]

%\setmonofont[
%  Path=./fonts,
%  Extension=.ttf,
%  Contextuals=NoAlternate,
%  UprightFont=*-regular,
%  BoldFont=*-bold,
%  ItalicFont=*-oblique,
%  SlantedFont=*-oblique,
%  Scale=0.87,
%]{iosevka-slab}

\def \MakeFrameTitles{
  \frametitle<1>{\FrameTitle}%
  \frametitle<2>{\FrameTitle\,--\,powtórzenie}%
}

\def \beginFrame#1{\begin{frame}<1>[fragile,label=#1]\MakeFrameTitles}
\def \beginMinted{\begin{minted}[tabsize=3,obeytabs]{go}}

\def \le {\leqslant}
\def \ge {\geqslant}
